\subsection{Formula}

Let $D$ be a district polygon and let MBC$(D)$ be the smallest circle that
contains $D$, with radius $r(D)$.
The Reock score is:
\begin{equation}
  \text{Reock}(D) = \frac{A(D)}{\pi r(D)^2},
  \label{eq:reock}
\end{equation}
where $A(D)$ is the area of $D$.

\begin{proposition}
  Reock$(D) \in (0, 1]$ for any non-degenerate polygon $D$,
  with equality if and only if $D$ is a disk.
\end{proposition}

\begin{proof}
  By definition, $D \subseteq \text{MBC}(D)$, so $A(D) \leq \pi r(D)^2$.
  Thus Reock$(D) \leq 1$.
  Equality holds iff $D = \text{MBC}(D)$ (as a set), i.e., iff $D$ is a disk.
  Reock$(D) > 0$ for any non-degenerate polygon since $A(D) > 0$. \qed
\end{proof}

\subsection{Canonical Values}

\begin{table}[ht]
\centering
\caption{Reock score for canonical shapes.}
\label{tab:reock_canonical}
\begin{tabular}{lcc}
\toprule
Shape & Reock & Approximate value \\
\midrule
Disk (circle) & $1$ & $1.000$ \\
Square of side $a$ & $2/\pi$ & $0.637$ \\
Regular hexagon & $3\sqrt{3}/(2\pi)$ & $0.827$ \\
$1 \times 3$ rectangle & $4/(5\pi)$ & $0.255$ \\
Equilateral triangle & $3\sqrt{3}/(4\pi)$ & $0.413$ \\
\bottomrule
\end{tabular}
\end{table}

\subsection{Properties}

\begin{proposition}
  Reock is scale-invariant, rotation-invariant, and translation-invariant.
\end{proposition}

\begin{proof}
  Scaling by $\lambda$: $A \mapsto \lambda^2 A$ and $r \mapsto \lambda r$,
  so Reock $= \lambda^2 A / (\pi \lambda^2 r^2) = A/(\pi r^2)$.
  Rotation and translation do not change areas or distances. \qed
\end{proof}

\begin{proposition}
  Reock$(D) \geq $ PP$(D)$ for any convex polygon $D$.
\end{proposition}

This follows because the MBC is contained within the isoperimetric circle
for any convex shape; the proof is omitted as it requires multiline
geometric argument beyond the scope of this paper.
